% =====================================================================
% tinjauan-pustaka.tex
% Tinjauan Pustaka Sistematis:
% YOLO, RGB, RGB+Depth, dan YOLO+RGB-D untuk Deteksi Objek (2019-2026)
%
% Kompilasi:
%   pdflatex tinjauan-pustaka
%   bibtex   tinjauan-pustaka
%   pdflatex tinjauan-pustaka
%   pdflatex tinjauan-pustaka
% =====================================================================
\documentclass[11pt,a4paper]{article}

\usepackage[utf8]{inputenc}
\usepackage[T1]{fontenc}
\usepackage[english,bahasa]{babel}
\usepackage{amsmath,amssymb}
\usepackage{graphicx}
\usepackage{booktabs}
\usepackage{longtable}
\usepackage{array}
\usepackage{multirow}
\usepackage[hidelinks]{hyperref}
\usepackage{geometry}
\geometry{margin=2.5cm}
\usepackage{cite}
\usepackage{microtype}

\title{\textbf{Tinjauan Pustaka Sistematis atas YOLO, Citra RGB, \\
Fusi RGB--Depth, dan Integrasi YOLO+RGB-D untuk \\
Deteksi Objek dan Persepsi Visual (2019--2026)}}
\author{Disusun secara autonomus\\ \small Research Pipeline --- Asisten Dosen}
\date{\today}

\begin{document}
\maketitle

\begin{abstract}
\noindent
Tinjauan pustaka ini merangkum perkembangan riset deteksi objek dan persepsi
visual dalam rentang 2019--2026 dengan empat poros utama: (i) evolusi keluarga
detektor satu-tahap \emph{You Only Look Once} (YOLO); (ii) deteksi objek berbasis
citra warna (RGB); (iii) pemanfaatan modalitas kedalaman melalui fusi RGB--Depth
(RGB-D); dan (iv) integrasi YOLO dengan data RGB-D. Sebanyak lebih dari 100 karya
ilmiah dianalisis dan dikelompokkan ke dalam sub-tema: fondasi detektor RGB,
survei YOLO, \emph{RGB-D salient object detection}, segmentasi semantik RGB-D,
estimasi kedalaman, estimasi pose 6D, deteksi \emph{grasp} robotik, deteksi 3D
fusi LiDAR--kamera, deteksi pedestrian multispektral, \emph{RGB-D SLAM}, serta
aplikasi lintas domain (pertanian, medis, industri, dan \emph{remote sensing}).
Temuan utama menunjukkan bahwa penambahan kanal kedalaman secara konsisten
memperbaiki ketahanan deteksi pada kondisi oklusi, pencahayaan rendah, dan
kebutuhan lokalisasi 3D, sementara YOLO tetap menjadi tulang punggung inferensi
waktu-nyata. Tantangan yang tersisa mencakup penyelarasan (alignment) antar
modalitas, keterbatasan sensor kedalaman pada jarak dekat dan cahaya matahari
kuat, serta kelangkaan dataset RGB-D beranotasi berskala besar.
\end{abstract}

\tableofcontents
\newpage

% =====================================================================
\section{Pendahuluan}
% =====================================================================
Deteksi objek adalah salah satu masalah fundamental dalam visi komputer yang
telah berkembang pesat selama dua dekade terakhir~\cite{zou2023objectdetectionsurvey}.
Sejak diperkenalkannya arsitektur \emph{region-based} seperti R-CNN dan
turunannya~\cite{girshick2014rcnn,girshick2015fastrcnn,ren2017fasterrcnn}, riset
bergeser menuju detektor satu-tahap yang mengutamakan kecepatan inferensi, dengan
SSD~\cite{liu2016ssd} dan YOLO~\cite{redmon2016yolo} sebagai representasi awal.
Kebutuhan aplikasi waktu-nyata --- kendaraan otonom, robotika, pertanian presisi,
serta diagnosis medis --- mendorong pengembangan detektor yang tidak hanya akurat
tetapi juga efisien~\cite{terven2023yolo,jiang2022yoloreview}.

Sebagian besar detektor beroperasi pada citra warna (RGB) dua dimensi. Namun,
citra RGB kehilangan informasi geometri sehingga rapuh terhadap oklusi,
tumpang tindih objek, dan pencahayaan yang buruk. Kehadiran sensor kedalaman
berbiaya rendah (misalnya Microsoft Kinect dan Intel RealSense) memungkinkan
akuisisi data RGB-D yang menyandingkan tekstur warna dengan geometri
\emph{per-pixel}~\cite{silberman2012nyu,song2015sunrgbd}. Fusi RGB dan kedalaman
terbukti memperkaya representasi fitur pada berbagai tugas: deteksi objek
\emph{salient}~\cite{zhou2021rgbdsurvey}, segmentasi semantik dalam
ruangan~\cite{hazirbas2016fusenet}, estimasi pose 6D~\cite{wang2019densefusion},
serta manipulasi robotik~\cite{morrison2018ggcnn}.

Tinjauan ini secara khusus menautkan kedua arus tersebut: bagaimana keluarga
YOLO yang ringan dan cepat dapat diperluas dengan modalitas kedalaman untuk
menghasilkan sistem persepsi yang lebih tangguh (\emph{YOLO+RGB-D}). Rentang
studi dibatasi pada 2019--2026 untuk menangkap tren mutakhir, dengan tetap
menyertakan karya-karya fondasi (2014--2018) sebagai konteks historis.

\subsection{Metodologi Penelusuran}
Literatur dikumpulkan melalui penelusuran bertahap pada mesin pencari akademik
dan basis data (arXiv, IEEE Xplore, ScienceDirect, SpringerLink, MDPI, dan PLOS)
menggunakan kombinasi kata kunci: \emph{YOLO}, \emph{object detection},
\emph{RGB-D}, \emph{depth fusion}, \emph{cross-modal}, \emph{salient object
detection}, \emph{semantic segmentation}, \emph{6D pose}, \emph{robotic
grasping}, \emph{3D detection}, dan turunannya. Kriteria inklusi meliputi
relevansi tema, kejelasan metrik, dan cakupan tahun 2019--2026 (dengan
pengecualian untuk karya fondasi). Total lebih dari 100 referensi disaring dan
disintesis ke dalam taksonomi tematik yang disajikan pada
Tabel~\ref{tab:taksonomi}.

% =====================================================================
\section{Evolusi Keluarga YOLO}
% =====================================================================
Sejak diperkenalkan oleh Redmon dkk.~\cite{redmon2016yolo}, YOLO merumuskan
deteksi objek sebagai satu masalah regresi tunggal yang memprediksi kotak
pembatas dan probabilitas kelas secara langsung dari piksel. YOLOv2/YOLO9000
memperkenalkan \emph{anchor box} dan pelatihan multi-skala~\cite{redmon2017yolo9000},
sedangkan YOLOv3 mengadopsi prediksi multi-skala dan \emph{backbone} Darknet-53
untuk objek kecil~\cite{redmon2018yolov3}. YOLOv4 mengonsolidasikan berbagai
teknik \emph{bag-of-freebies} dan \emph{bag-of-specials}~\cite{bochkovskiy2020yolov4},
diikuti varian industri PP-YOLO~\cite{long2020ppyolo} dan YOLOX yang beralih ke
desain \emph{anchor-free} dan \emph{decoupled head}~\cite{ge2021yolox}.

Generasi berikutnya menekankan optimasi proses pelatihan dan efisiensi:
YOLOv6 dioptimalkan untuk aplikasi industri~\cite{li2022yolov6}, YOLOv7
memperkenalkan \emph{trainable bag-of-freebies} dan \emph{E-ELAN}~\cite{wang2023yolov7},
YOLOv9 menghadirkan \emph{Programmable Gradient Information} dan
\emph{GELAN}~\cite{wang2024yolov9}, sedangkan YOLOv10 menghapus kebutuhan
\emph{Non-Maximum Suppression} melalui pelatihan ujung-ke-ujung yang
konsisten~\cite{wang2024yolov10}. YOLOv11 menambahkan modul C3K2 dan C2PSA
berbasis atensi~\cite{khanam2024yolov11}. Sejumlah survei komprehensif memetakan
lintasan ini secara sistematis~\cite{terven2023yolo,hussain2023yolo,diwan2023yolo,
sapkota2024yoloagri,ali2024yoloreview,vijayakumar2024yolo,alif2024yoloevolution,
sohan2024yolov8review}. Ringkasan evolusi disajikan pada Tabel~\ref{tab:yolo}.

\begin{table}[htbp]
\centering
\caption{Ringkasan tonggak evolusi keluarga YOLO.}
\label{tab:yolo}
\small
\begin{tabular}{@{}llp{7.6cm}@{}}
\toprule
\textbf{Versi} & \textbf{Tahun} & \textbf{Kontribusi utama} \\
\midrule
YOLOv1~\cite{redmon2016yolo}      & 2016 & Deteksi sebagai regresi tunggal, sangat cepat \\
YOLOv2~\cite{redmon2017yolo9000}  & 2017 & Anchor box, batch-norm, pelatihan multi-skala \\
YOLOv3~\cite{redmon2018yolov3}    & 2018 & Prediksi multi-skala, Darknet-53 \\
YOLOv4~\cite{bochkovskiy2020yolov4} & 2020 & Bag-of-freebies/specials, CSPDarknet \\
YOLOX~\cite{ge2021yolox}          & 2021 & Anchor-free, decoupled head, SimOTA \\
YOLOv6~\cite{li2022yolov6}        & 2022 & Desain industri, reparameterisasi \\
YOLOv7~\cite{wang2023yolov7}      & 2022 & Trainable bag-of-freebies, E-ELAN \\
YOLOv9~\cite{wang2024yolov9}      & 2024 & PGI, GELAN, mitigasi kehilangan informasi \\
YOLOv10~\cite{wang2024yolov10}    & 2024 & Bebas-NMS, pelatihan ujung-ke-ujung \\
YOLOv11~\cite{khanam2024yolov11}  & 2024 & Modul C3K2, C2PSA berbasis atensi \\
\bottomrule
\end{tabular}
\end{table}

% =====================================================================
\section{Deteksi Objek Berbasis RGB: Fondasi dan Transformer}
% =====================================================================
Selain YOLO, ekosistem deteksi RGB mencakup detektor dua-tahap dan satu-tahap
lain yang menjadi pembanding penting. RetinaNet memperkenalkan \emph{focal loss}
untuk menangani ketidakseimbangan kelas~\cite{lin2017focal}, sementara
\emph{Feature Pyramid Network} menjadi standar de facto untuk fitur
multi-skala~\cite{lin2017fpn}. Pendekatan \emph{anchor-free} seperti
FCOS~\cite{tian2019fcos} dan CenterNet~\cite{zhou2019objects} menyederhanakan
alur deteksi. EfficientDet mengusung \emph{BiFPN} dan penskalaan
majemuk~\cite{tan2020efficientdet}, sedangkan Mask R-CNN memperluas deteksi ke
segmentasi instans~\cite{he2017maskrcnn}.

Transformer membawa paradigma baru: DETR merumuskan deteksi sebagai prediksi
himpunan ujung-ke-ujung~\cite{carion2020detr}, disempurnakan oleh Deformable
DETR untuk konvergensi lebih cepat~\cite{zhu2021deformabledetr}. \emph{Vision
Transformer}~\cite{dosovitskiy2021vit} dan \emph{Swin Transformer} dengan jendela
bergeser~\cite{liu2021swin} menjadi \emph{backbone} tugas prediksi padat.
\emph{Backbone} residual~\cite{he2016resnet} dan modul atensi
CBAM~\cite{woo2018cbam} tetap menjadi komponen umum. Dataset COCO~\cite{lin2014coco}
menjadi tolok ukur standar untuk seluruh keluarga detektor ini.

% =====================================================================
\section{Modalitas Kedalaman dan Estimasi Depth}
% =====================================================================
Data kedalaman dapat diperoleh langsung dari sensor RGB-D atau diestimasi dari
citra. Estimasi kedalaman monokular berkembang dari regresi multi-skala
awal~\cite{eigen2014depth} menuju pendekatan swa-penyeliaan yang memanfaatkan
konsistensi fotometrik stereo maupun video~\cite{godard2017monodepth,
godard2019monodepth2,guizilini2020packnet}. Metode berbasis \emph{adaptive bins}
~\cite{bhat2021adabins} dan transformer untuk prediksi
padat~\cite{ranftl2021dpt} meningkatkan akurasi, sementara MiDaS mencapai
generalisasi lintas-dataset (\emph{zero-shot})~\cite{ranftl2022midas} dan Depth
Anything memanfaatkan data tak-berlabel berskala besar~\cite{yang2024depthanything}.
Kasus dalam-ruangan ditangani secara khusus oleh
MonoIndoor~\cite{ji2021monoindoor}, dan tinjauan menyeluruh atas estimasi
kedalaman berbasis pembelajaran mendalam disajikan oleh Ming
dkk.~\cite{ming2021depthsurvey}. Kualitas peta kedalaman --- baik terukur
maupun terestimasi --- menjadi penentu penting keberhasilan fusi RGB-D.

% =====================================================================
\section{RGB-D Salient Object Detection (SOD)}
% =====================================================================
\emph{Salient object detection} berbasis RGB-D bertujuan menonjolkan objek yang
paling menarik perhatian dengan bantuan isyarat geometri. DMRA memanfaatkan
atensi rekuren multi-skala yang diinduksi kedalaman~\cite{piao2019dmra},
sedangkan BBS-Net mengusung strategi \emph{bifurcated backbone}~\cite{fan2020bbsnet}.
Fan dkk. menata ulang tolok ukur dan model bidang ini secara
komprehensif~\cite{fan2020d3net}. Sejumlah strategi fusi lintas-modal
dikembangkan: JL-DCF dengan pembelajaran bersama~\cite{fu2020jldcf},
\emph{self-mutual attention}~\cite{liu2020s2ma}, penyaringan dinamis
hierarkis~\cite{pang2020hdfnet}, serta pemodelan ketidakpastian melalui
\emph{variational autoencoder}~\cite{zhang2020ucnet}.

Gelombang transformer menandai pergeseran arsitektur: \emph{Visual Saliency
Transformer}~\cite{liu2021vst}, SwinNet yang menyadari tepi untuk RGB-D dan
RGB-T~\cite{liu2021swinnet}, serta TriTransNet dengan
\emph{triplet transformer}~\cite{liu2021tritransnet}. Pendekatan atensi
sensitif-kedalaman dan fusi multi-modal otomatis diusulkan
DSA2F~\cite{sun2021dsa2f}. Karya termutakhir menekankan interaksi dan penghalusan
lintas-modal~\cite{cong2022cirnet}, penyaringan dinamis
\emph{criss-cross}~\cite{zhang2022c2dfnet}, fusi adaptif lintas-skala
~\cite{zhou2022ccafnet}, interaksi hierarkis~\cite{chen2023cmhi}, serta kesadaran
kedalaman hierarkis~\cite{wu2023hidanet}. Survei Zhou
dkk.~\cite{zhou2021rgbdsurvey} merangkum tren, dataset, dan metrik evaluasi
bidang ini.

% =====================================================================
\section{Segmentasi Semantik RGB-D}
% =====================================================================
Segmentasi semantik dalam-ruangan sangat diuntungkan oleh kedalaman untuk
memisahkan objek berdekatan dengan tekstur serupa. FuseNet mengusung fusi
berbasis enkoder ganda~\cite{hazirbas2016fusenet}, RedNet memanfaatkan struktur
enkoder-dekoder residual~\cite{jiang2018rednet}, dan RDFNet menyusun fusi fitur
residual bertingkat~\cite{park2017rdfnet}. ACNet menimbang fitur antar modalitas
melalui atensi~\cite{hu2019acnet}, sedangkan SA-Gate mengatur propagasi fitur
lintas-modalitas dua arah~\cite{chen2020sagate}. ESANet menargetkan efisiensi
untuk analisis \emph{scene} dalam-ruangan pada platform waktu-nyata
~\cite{seichter2021esanet}, dan ShapeConv menyuntikkan kesadaran bentuk ke dalam
lapisan konvolusi~\cite{cao2021shapeconv}.

Pendekatan berbasis transformer memperluas cakupan: TokenFusion memfusikan token
multi-modal pada ViT~\cite{wang2022tokenfusion}, CMX menyediakan kerangka fusi
RGB-X yang umum~\cite{zhang2023cmx}, PGDENet menerapkan fusi terpandu progresif
dan penguatan kedalaman~\cite{zhou2022pgdenet}, dan DFormer memikirkan ulang
pembelajaran representasi RGB-D untuk segmentasi~\cite{yin2024dformer}. Dataset
NYU Depth v2~\cite{silberman2012nyu}, SUN RGB-D~\cite{song2015sunrgbd}, dan
ScanNet~\cite{dai2017scannet} menjadi tolok ukur utama.

% =====================================================================
\section{Fusi RGB+Depth untuk Deteksi Objek}
% =====================================================================
Fusi RGB dan kedalaman untuk deteksi objek umumnya dikategorikan menjadi fusi
awal (\emph{early}/input-level), fusi menengah (\emph{intermediate}/feature-level),
dan fusi akhir (\emph{late}/decision-level)~\cite{ramachandram2017multimodalreview,
feng2021multimodalsurvey}. Ophoff dkk. menyelidiki strategi fusi RGB+Depth untuk
deteksi waktu-nyata dan menemukan bahwa fusi tingkat-fitur menengah memberikan
keseimbangan terbaik antara akurasi dan kecepatan~\cite{ophoff2019multimodal}.
Farahnakian dan Heikkonen mengusulkan strategi fusi berbasis atensi untuk
mengintegrasikan RGB dan kedalaman secara efektif~\cite{farahnakian2021fusion}.
Survei atas dataset RGB-D~\cite{lopes2022rgbddatasets} menyoroti keragaman sensor
dan tantangan penyelarasan. Tema berulang dalam literatur ini adalah bahwa
kualitas dan penyelarasan peta kedalaman menentukan seberapa besar keuntungan
fusi dapat direalisasikan.

% =====================================================================
\section{Integrasi YOLO dengan RGB-D}
% =====================================================================
Integrasi YOLO dengan RGB-D menempuh beberapa strategi. Pendekatan paling
langsung memperluas kanal masukan: \emph{Expandable YOLO} menambahkan kanal
kedalaman ke tiga kanal RGB YOLOv3 sehingga jaringan mengekstraksi fitur warna
dan geometri secara simultan untuk menghasilkan kotak pembatas
3D~\cite{takahashi2020expandableyolo}. Pendekatan lain memakai YOLO sebagai
detektor 2D lalu memproyeksikan hasilnya ke peta kedalaman teraligned untuk
lokalisasi 3D, sebagaimana pada sistem navigasi robot~\cite{xu2024onboard} dan
pengukuran jarak objek dalam-ruangan~\cite{chen2023depthyolo}. FusionVision
memadukan deteksi YOLO dengan segmentasi cepat untuk rekonstruksi dan segmentasi
objek 3D dari kamera RGB-D~\cite{elamraoui2024fusionvision}.

Pada manipulasi robotik, Tian dkk. memfusikan fitur RGB dan kedalaman di dalam
kepala prediksi berbasis YOLO untuk deteksi \emph{grasp}~\cite{tian2023grasp},
sedangkan YOLOv8-URE menggabungkan visi 2D dan awan titik untuk penggenggaman
pada skenario objek bertumpuk~\cite{yang2025yolov8ure}. Di sektor pertanian,
robot pemetik labu memanfaatkan kamera RGB-D dan model deteksi berbasis
YOLO~\cite{ito2024pumpkin}. Studi-studi ini secara konsisten melaporkan bahwa
penambahan kedalaman meningkatkan ketahanan terhadap oklusi serta memungkinkan
estimasi posisi 3D yang diperlukan untuk aksi robotik.

% =====================================================================
\section{Estimasi Pose 6D dan Deteksi Grasp Berbasis RGB-D}
% =====================================================================
Estimasi pose 6D menuntut lokalisasi translasi dan rotasi objek secara akurat.
PoseCNN memperkenalkan regresi pose langsung dari CNN~\cite{xiang2018posecnn},
sedangkan DenseFusion memfusikan fitur RGB dan awan titik secara padat dan
iteratif~\cite{wang2019densefusion}. Pendekatan berbasis \emph{voting}
titik-kunci seperti PVN3D~\cite{he2020pvn3d} dan fusi dua-arah aliran-penuh
FFB6D~\cite{he2021ffb6d} meningkatkan ketahanan pada \emph{scene} berantakan.
G2L-Net menawarkan inferensi waktu-nyata~\cite{chen2020g2lnet}, dan FoundationPose
menyatukan estimasi serta pelacakan objek baru~\cite{wen2024foundationpose}.
Survei komprehensif merangkum kemajuan bidang ini~\cite{hoque2021posesurvey}.

Deteksi \emph{grasp} robotik berevolusi dari klasifikasi kandidat
\emph{grasp}~\cite{lenz2015grasp} menuju sintesis generatif waktu-nyata: GG-CNN
memprediksi kualitas \emph{grasp} per-piksel untuk kendali
lingkar-tertutup~\cite{morrison2018ggcnn}, sedangkan GR-ConvNet dan
penerusnya menghasilkan \emph{grasp} antipodal dari masukan RGB-D dengan akurasi
tinggi pada dataset Cornell dan Jacquard~\cite{kumra2020grconvnet,
kumra2022grconvnetv2,depierre2018jacquard}. Jaringan fusi lintas-modal bilateral
lebih jauh menyelaraskan isyarat warna dan geometri~\cite{zhang2023bcmfnet},
sementara GraspNet-1Billion menyediakan tolok ukur berskala
besar~\cite{fang2020graspnet}.

% =====================================================================
\section{Deteksi 3D: Fusi LiDAR/Awan Titik dan Kamera}
% =====================================================================
Untuk kendaraan otonom, kedalaman kerap berasal dari LiDAR. Pendekatan berbasis
voxel~\cite{zhou2018voxelnet} dan \emph{pillar}~\cite{lang2019pointpillars}
memproses awan titik secara efisien, sedangkan PointRCNN membangkitkan proposal
langsung dari titik~\cite{shi2019pointrcnn}. Frustum PointNets memadukan deteksi
2D RGB dengan awan titik RGB-D~\cite{qi2018frustum}, mengandalkan
\emph{backbone} PointNet~\cite{qi2017pointnet}. Fusi kamera--LiDAR ditempuh
melalui agregasi multi-tampilan~\cite{chen2017mv3d,ku2018avod}, pewarnaan titik
sekuensial~\cite{vora2020pointpainting}, fusi fitur lintas-tampilan
~\cite{yoo20203dcvf}, hingga pendekatan transformer~\cite{bai2022transfusion} dan
representasi \emph{bird's-eye view} terpadu~\cite{liu2023bevfusion}. Alternatif
\emph{pseudo-LiDAR} merekonstruksi geometri dari citra stereo/monokular
~\cite{wang2019pseudolidar,you2020pseudolidarpp}. Dataset
KITTI~\cite{geiger2012kitti} dan nuScenes~\cite{caesar2020nuscenes} menjadi tolok
ukur utama, dengan survei yang merangkum metode dan
tantangannya~\cite{arnold2019survey3d,feng2021multimodalsurvey}.

% =====================================================================
\section{Deteksi Pedestrian Multispektral (RGB-Thermal)}
% =====================================================================
Meski memakai termal alih-alih kedalaman, deteksi pedestrian multispektral
berbagi prinsip fusi lintas-modal yang sama. Tolok ukur KAIST menjadi rujukan
standar~\cite{hwang2015kaist}. Li dkk. mengusulkan Faster R-CNN yang sadar
pencahayaan~\cite{li2019illumination}, sedangkan MBNet menangani ketidakseimbangan
modalitas~\cite{zhou2020mbnet}. Strategi fusi berbasis atensi terpandu
~\cite{zhang2021gaff}, blok fuse-and-refine siklik~\cite{zhang2020cfr}, serta
pembelajaran lintas-modal terpandu-ketidakpastian~\cite{kim2022cmpd} meningkatkan
ketahanan pada kondisi malam dan cahaya rendah. Pelajaran dari domain ini ---
penanganan ketidakseimbangan modalitas dan penyelarasan lemah --- berlaku pula
untuk fusi RGB-D.

% =====================================================================
\section{RGB-D SLAM Dinamis dan Semantik}
% =====================================================================
Pada robotika bergerak, RGB-D dimanfaatkan untuk pemetaan dan lokalisasi
simultan. ORB-SLAM2 menyediakan kerangka SLAM untuk kamera monokular, stereo, dan
RGB-D~\cite{murartal2017orbslam2}. Lingkungan dinamis ditangani dengan menggabungkan
segmentasi/deteksi: DynaSLAM memisahkan objek bergerak dan melakukan
\emph{inpainting}~\cite{bescos2018dynaslam}, DS-SLAM memadukan SLAM semantik dengan
penolakan fitur dinamis~\cite{yu2018dsslam}, dan CFP-SLAM memakai probabilitas
kasar-ke-halus secara waktu-nyata~\cite{hu2022cfpslam}. Sejumlah sistem
memanfaatkan detektor seperti YOLO atau Mask R-CNN untuk menandai objek dinamis,
dengan trade-off antara akurasi dan kecepatan yang perlu
diseimbangkan~\cite{soares2019visualslam}.

% =====================================================================
\section{Aplikasi Lintas Domain}
% =====================================================================
\subsection{Pertanian dan Panen Robotik}
YOLO banyak diadopsi untuk deteksi buah waktu-nyata: MangoYOLO untuk estimasi
beban panen~\cite{koirala2019deepfruit}, YOLOv3 termodifikasi untuk deteksi apel
pada berbagai tahap pertumbuhan~\cite{tian2019appleyolo}, dan YOLOv4 terpangkas
untuk deteksi bunga apel~\cite{wu2020flower}. Kamera RGB-D memungkinkan lokalisasi
3D buah untuk penggenggaman: deteksi apel berbasis fitur RGB dan
kedalaman~\cite{fu2020apple}, lokalisasi 3D melalui segmentasi
instans~\cite{genemola2020fruit3d}, deteksi dan visualisasi 3D di kebun
apel~\cite{kang2020strawberry}, sistem panen selada~\cite{birrell2020lettuce},
robot pemetik buah otomatis~\cite{onishi2019harvest}, dan robot pemetik labu
berbasis RGB-D+YOLO~\cite{ito2024pumpkin}. Survei aplikasi YOLO di domain
pertanian merangkum tren ini~\cite{sapkota2024yoloagri}.

\subsection{Medis}
YOLO telah diterapkan luas pada pencitraan medis, sebagaimana dirangkum oleh
tinjauan sistematis~\cite{qureshi2024medyolo}. Contoh spesifik mencakup diagnosis
berbasis citra dada~\cite{alantari2020covid}, deteksi dan klasifikasi lesi
payudara~\cite{baccouche2021breast,mohiyuddin2022breast}, serta deteksi polip
pada kolonoskopi waktu-nyata~\cite{wan2023polyp}.

\subsection{Industri dan Inspeksi}
Untuk inspeksi permukaan, EFC-YOLO mendeteksi cacat pada pelat
baja~\cite{yang2023efcyolo} dan PCB-YOLO meningkatkan deteksi cacat papan sirkuit
melalui atensi koordinat dan fusi multi-skala~\cite{tang2024pcbyolo}. Aplikasi
keselamatan mencakup deteksi pemakaian helm~\cite{zhou2021helmet}, dengan tinjauan
lebih luas atas deteksi cacat berbasis pembelajaran mendalam~\cite{jha2023defectreview}
dan peran YOLO dalam manufaktur digital~\cite{hussain2023yolo}.

\subsection{Remote Sensing dan UAV}
Deteksi objek kecil pada citra udara ditangani TPH-YOLOv5 dengan kepala prediksi
transformer~\cite{zhu2021tphyolov5}, UAV-YOLOv8 untuk skenario fotografi
udara~\cite{wang2023uavyolo}, YOLT untuk citra satelit~\cite{vanetten2018yolt},
serta CNN hierarkis untuk penginderaan jauh resolusi
tinggi~\cite{zhang2019remotesensing}.

% =====================================================================
\section{Taksonomi, Sintesis, dan Diskusi}
% =====================================================================
Tabel~\ref{tab:taksonomi} merangkum taksonomi tematik yang digunakan dalam
tinjauan ini beserta karya-karya representatifnya.

\begin{longtable}{@{}p{3.2cm}p{4.0cm}p{6.3cm}@{}}
\caption{Taksonomi tematik literatur yang ditinjau (2019--2026 dan karya fondasi).}
\label{tab:taksonomi}\\
\toprule
\textbf{Tema} & \textbf{Fokus} & \textbf{Karya representatif} \\
\midrule
\endfirsthead
\toprule
\textbf{Tema} & \textbf{Fokus} & \textbf{Karya representatif} \\
\midrule
\endhead
\bottomrule
\endfoot
Evolusi YOLO & Detektor satu-tahap waktu-nyata &
\cite{redmon2016yolo,redmon2017yolo9000,redmon2018yolov3,bochkovskiy2020yolov4,
ge2021yolox,li2022yolov6,wang2023yolov7,wang2024yolov9,wang2024yolov10,
khanam2024yolov11,long2020ppyolo} \\
Survei YOLO & Pemetaan sistematis &
\cite{terven2023yolo,hussain2023yolo,jiang2022yoloreview,diwan2023yolo,
sapkota2024yoloagri,ali2024yoloreview,vijayakumar2024yolo,alif2024yoloevolution,
sohan2024yolov8review} \\
Fondasi RGB & Detektor \& transformer &
\cite{girshick2014rcnn,girshick2015fastrcnn,ren2017fasterrcnn,liu2016ssd,
lin2017focal,he2017maskrcnn,lin2017fpn,tian2019fcos,zhou2019objects,
tan2020efficientdet,carion2020detr,zhu2021deformabledetr,dosovitskiy2021vit,
liu2021swin,he2016resnet,woo2018cbam} \\
Estimasi kedalaman & Monokular \& swa-penyeliaan &
\cite{eigen2014depth,godard2017monodepth,godard2019monodepth2,guizilini2020packnet,
bhat2021adabins,ranftl2021dpt,ranftl2022midas,yang2024depthanything,
ji2021monoindoor,ming2021depthsurvey} \\
RGB-D SOD & Fusi lintas-modal saliency &
\cite{piao2019dmra,fan2020bbsnet,fan2020d3net,fu2020jldcf,liu2020s2ma,
pang2020hdfnet,zhang2020ucnet,liu2021vst,liu2021swinnet,liu2021tritransnet,
sun2021dsa2f,cong2022cirnet,zhang2022c2dfnet,zhou2022ccafnet,chen2023cmhi,
wu2023hidanet,zhou2021rgbdsurvey} \\
Segmentasi RGB-D & Scene parsing dalam-ruangan &
\cite{hazirbas2016fusenet,jiang2018rednet,park2017rdfnet,hu2019acnet,
chen2020sagate,seichter2021esanet,cao2021shapeconv,wang2022tokenfusion,
zhang2023cmx,zhou2022pgdenet,yin2024dformer} \\
Fusi RGB+Depth & Deteksi objek multimodal &
\cite{ophoff2019multimodal,farahnakian2021fusion,ramachandram2017multimodalreview,
feng2021multimodalsurvey,lopes2022rgbddatasets} \\
YOLO+RGB-D & Integrasi kedalaman ke YOLO &
\cite{takahashi2020expandableyolo,elamraoui2024fusionvision,tian2023grasp,
yang2025yolov8ure,ito2024pumpkin,xu2024onboard,chen2023depthyolo} \\
Pose 6D \& grasp & Manipulasi robotik &
\cite{xiang2018posecnn,wang2019densefusion,he2020pvn3d,he2021ffb6d,chen2020g2lnet,
wen2024foundationpose,hoque2021posesurvey,lenz2015grasp,morrison2018ggcnn,
kumra2020grconvnet,kumra2022grconvnetv2,zhang2023bcmfnet,fang2020graspnet,
depierre2018jacquard} \\
Deteksi 3D & Fusi LiDAR--kamera &
\cite{zhou2018voxelnet,lang2019pointpillars,shi2019pointrcnn,qi2018frustum,
qi2017pointnet,chen2017mv3d,ku2018avod,vora2020pointpainting,yoo20203dcvf,
bai2022transfusion,liu2023bevfusion,wang2019pseudolidar,you2020pseudolidarpp,
arnold2019survey3d} \\
Pedestrian RGB-T & Fusi multispektral &
\cite{hwang2015kaist,li2019illumination,zhou2020mbnet,zhang2021gaff,zhang2020cfr,
kim2022cmpd} \\
RGB-D SLAM & Lingkungan dinamis &
\cite{murartal2017orbslam2,bescos2018dynaslam,yu2018dsslam,hu2022cfpslam,
soares2019visualslam} \\
Aplikasi & Pertanian, medis, industri, RS &
\cite{koirala2019deepfruit,tian2019appleyolo,wu2020flower,fu2020apple,
genemola2020fruit3d,kang2020strawberry,birrell2020lettuce,onishi2019harvest,
qureshi2024medyolo,alantari2020covid,baccouche2021breast,mohiyuddin2022breast,
wan2023polyp,yang2023efcyolo,tang2024pcbyolo,jha2023defectreview,zhou2021helmet,
zhu2021tphyolov5,wang2023uavyolo,vanetten2018yolt,zhang2019remotesensing} \\
Dataset & Tolok ukur RGB-D/deteksi &
\cite{silberman2012nyu,song2015sunrgbd,dai2017scannet,geiger2012kitti,
caesar2020nuscenes,lin2014coco,zou2023objectdetectionsurvey} \\
\end{longtable}

\subsection{Diskusi Lintas-Tema}
Sintesis literatur mengungkap beberapa pola. \emph{Pertama}, terdapat konvergensi
menuju arsitektur hibrida CNN--transformer, baik pada detektor
RGB~\cite{liu2021swin,carion2020detr} maupun pada fusi
RGB-D~\cite{liu2021vst,yin2024dformer}. \emph{Kedua}, strategi fusi tingkat-fitur
menengah dengan mekanisme atensi lintas-modal secara konsisten mengungguli fusi
awal atau akhir~\cite{chen2020sagate,ophoff2019multimodal,cong2022cirnet}.
\emph{Ketiga}, YOLO tetap menjadi pilihan utama untuk komponen deteksi 2D dalam
pipeline RGB-D karena keseimbangan kecepatan--akurasinya, sehingga banyak sistem
memakainya sebagai tahap awal sebelum penalaran geometri
3D~\cite{takahashi2020expandableyolo,tian2023grasp,xu2024onboard,ito2024pumpkin}.

% =====================================================================
\section{Tantangan dan Arah Riset Masa Depan}
% =====================================================================
Beberapa tantangan terbuka teridentifikasi dari tinjauan ini:
\begin{itemize}
  \item \textbf{Kualitas dan penyelarasan kedalaman.} Sensor RGB-D mengalami
  degradasi pada jarak dekat serta di bawah cahaya matahari kuat, menurunkan
  kualitas peta kedalaman~\cite{ophoff2019multimodal,lopes2022rgbddatasets}.
  \item \textbf{Ketidakseimbangan dan penyelarasan lemah antar modalitas},
  sebagaimana dipelajari pada domain multispektral~\cite{zhou2020mbnet,zhang2021gaff}.
  \item \textbf{Kelangkaan dataset RGB-D beranotasi berskala besar} untuk deteksi
  objek waktu-nyata, dibanding melimpahnya dataset RGB~\cite{lin2014coco,
  lopes2022rgbddatasets}.
  \item \textbf{Efisiensi komputasi} fusi multi-modal pada perangkat tepi
  (\emph{edge}), yang menuntut model ringan~\cite{seichter2021esanet,li2022yolov6}.
  \item \textbf{Generalisasi lintas-domain}, di mana estimasi kedalaman fondasi
  seperti Depth Anything~\cite{yang2024depthanything} berpotensi menyediakan
  kedalaman \emph{pseudo} yang murah untuk memperkaya YOLO tanpa sensor khusus.
\end{itemize}
Arah menjanjikan mencakup penggabungan YOLO generasi baru (v9--v11) dengan
kedalaman terestimasi fondasi, kerangka fusi transformer yang hemat parameter,
serta strategi pelatihan yang tangguh terhadap kegagalan/kehilangan modalitas.

% =====================================================================
\section{Kesimpulan}
% =====================================================================
Tinjauan atas lebih dari 100 karya (2019--2026, dengan konteks fondasi) memetakan
empat poros riset yang saling terkait: evolusi YOLO, deteksi RGB, fusi RGB-D, dan
integrasi YOLO+RGB-D. Bukti empiris lintas domain --- pertanian, medis, industri,
robotika, dan kendaraan otonom --- menunjukkan bahwa modalitas kedalaman
melengkapi kelemahan citra RGB pada oklusi, pencahayaan buruk, dan kebutuhan
lokalisasi 3D, sementara YOLO menyediakan tulang punggung inferensi waktu-nyata.
Integrasi keduanya, baik melalui perluasan kanal masukan maupun pipeline
deteksi-lalu-proyeksi, terbukti meningkatkan ketahanan sistem persepsi. Tantangan
utama yang tersisa berkisar pada kualitas serta penyelarasan kedalaman, efisiensi
fusi, dan ketersediaan data --- yang menjadi agenda riset yang menjanjikan ke
depan.

% =====================================================================
\bibliographystyle{ieeetr}
\bibliography{references}

\end{document}
